\scanpage{5}
Descriptions conjecturales de motifs sur corps de nbs \deleted{$\unsure{F}$} $k$

1) Double \unsure{réseau|viseur} Betti-deRham

2) Double \unsure{réseau|viseur} Betti-Hodge

3) \deleted{Betti + \unsure{bigrad}} structure de Hodge ($k=\overline{\mathbf Q}$)

4) \deleted{Betti + \unsure{connexion} \unsure{réseau} Betti}

Betti -- (Gauss)Manin \qquad ($k=\overline{\mathbf Q}$)

5) \deleted{Tate} Betti-Tate

6) \deleted{De Rham-cristal} De Rham filtré cristallin

Descriptions conjecturales des groupes $p$-divisibles \inserted{mod isogénie} (\unsure{plus génér\textsuperscript{t}} de pseudo-motifs) sur un corps local [\deleted{\unsure{spectre}}\inserted{corps de fr.} d'un ann. de val. \unsure{discrets|discrète} complet à corps rés $p$-fait] :

3') structure de Hodge-\inserted{Tate; module $N$ de d.f. sur $\mathbf Q_p$, et bigrad. sur $N\otimes_{\mathbf Q_p}C$,} (mod ext. finie de $K$)

2') double \unsure{réseau|viseur} Tate-Hodge

5') Tate (module galoisien $\ell$-adique)

6') \deleted{module \unsure{Vect. de}} cristal filtré [vectoriel \deleted{\unsure{sur}} dim finie sur $K$, avec application $\sigma$-semi-linéaire, et filtration, graduation]

Par quoi remplacer $1',4'$ ??
\mspage{4}
